\chapter{IMPLEMENTATION}
\section{Pseudocode}
% or Algorithm for the proposed work
\begin{algorithm}[H]
\caption{Federated Ensemble Fraud Detection Training}
\begin{algorithmic}[1]
\State Load CSV dataset from \texttt{data/}
\State Detect fraud label column and validate binary labels
\State Split data into training and testing sets
\State Fit preprocessing pipeline on training data
\State Transform training and testing features
\State Partition transformed training data into $K$ clients
\State Initialize global Logistic Regression and MLP parameters
\For{each federated round $r = 1$ to $R$}
    \For{each selected client $k$}
        \State Send global LR and MLP parameters to client $k$
        \State Train LR and MLP locally on client data
        \State Return updated parameters and local sample count $n_k$
    \EndFor
    \State Aggregate LR and MLP parameters using weighted FedAvg
\EndFor
\For{each client $k$}
    \State Train local XGBoost model on client data
\EndFor
\State Compute LR, MLP, and averaged XGBoost probabilities
\State Combine probabilities using weighted ensemble
\State Tune threshold on validation data
\State Evaluate on test data
\State Save models, metrics, plots, and metadata
\end{algorithmic}
\end{algorithm}

\section{Implementation Details}  %Algorithm and how it is being implemented}
%Information about the implementation of Modules			
\subsection{Data Loading and Preprocessing}
The module \texttt{training/data\_loader.py} loads the first CSV file found in the \texttt{data/} directory. The preprocessing module validates the label, separates features from labels, imputes missing values, scales numeric features, and one-hot encodes categorical features when present.

\begin{lstlisting}[language=Python, caption={Dataset Preparation Entry Point}]
client_partitions, test, preprocessor, feature_names, dataset_path = prepare_data(config)
\end{lstlisting}

The fitted preprocessor is saved using \texttt{joblib}, ensuring that inference uses the same transformations as training.

\subsection{Federated Client Implementation}
Each client is implemented as a Flower-compatible NumPy client. It receives global parameters, loads them into local models, trains on local data, and returns updated parameters.

\begin{lstlisting}[language=Python, caption={Client Fit Method Summary}]
def fit(self, parameters, config):
    self.set_parameters(parameters)
    self.logistic.fit(X, y, epochs=epochs)
    losses = self.mlp.fit(X, y, epochs=epochs, batch_size=batch_size)
    return self.get_parameters(), len(y), metrics
\end{lstlisting}

Only model parameters and metrics are returned. Raw rows are not sent to the server.

\subsection{Logistic Regression}
The Logistic Regression model uses stochastic gradient updates through scikit-learn's \texttt{SGDClassifier}. Its parameters are:
\begin{equation}
\theta_{LR} = \{w, b\}
\end{equation}
where $w$ is the coefficient vector and $b$ is the intercept. The fraud probability is:
\begin{equation}
P(y=1|x) = \sigma(w^Tx + b)
\end{equation}

\subsection{MLP}
The MLP is implemented in PyTorch with configurable hidden layers. The current configuration uses two hidden layers with 64 and 32 neurons and dropout of 0.2. The forward pass is:
\begin{equation}
h_1 = ReLU(W_1x+b_1), \quad h_2 = ReLU(W_2h_1+b_2), \quad z = W_3h_2+b_3
\end{equation}
and the output probability is $\sigma(z)$.

\subsection{XGBoost}
Each client trains its own XGBoost model locally. Since tree structures cannot be averaged like neural network weights, the system averages XGBoost probability predictions rather than parameters.

\begin{lstlisting}[language=Python, caption={Averaging Client XGBoost Predictions}]
scores = [model.predict_proba(X) for model in xgb_models]
xgb_scores = np.mean(np.vstack(scores), axis=0)
\end{lstlisting}

\subsection{Ensemble Aggregation}
The final ensemble is implemented in \texttt{ensemble/aggregator.py}. The default configuration is:
\begin{equation}
s = 0.3s_{LR} + 0.3s_{MLP} + 0.4s_{XGB}
\end{equation}
The final decision is:
\begin{equation}
\hat{y} =
\begin{cases}
1, & s \geq \tau \\
0, & s < \tau
\end{cases}
\end{equation}
where $\tau$ is the tuned threshold.

\subsection{Checkpointing}
The following artifacts are saved:
\begin{itemize}
    \item \texttt{logistic\_regression.joblib}
    \item \texttt{mlp.pth}
    \item \texttt{xgboost\_client\_*.json}
    \item \texttt{preprocessor.joblib}
    \item \texttt{ensemble\_config.json}
    \item \texttt{model\_metadata.json}
    \item \texttt{metrics.json}
\end{itemize}

\subsection{Inference API}
The FastAPI layer loads all checkpoints at startup and exposes \texttt{/predict}. The endpoint preprocesses incoming feature dictionaries, computes model scores, applies the ensemble, and returns the fraud probability and decision.
